% from ega3-1-fr.tex
\oldpage[0\textsubscript{III}]{35}
\[
  \label{0.11.4.6.1-fr}
  'E_2^{pq}=\mathrm H^p(\mathrm R^qT(K^\bullet,K'^\bullet))
  \tag{11.4.6.1}
\]
\[
  \label{0.11.4.6.2-fr}
  ''E_2^{pq}=\sum_{q'+q''=q}
  \mathrm R^pT(\mathrm H^{q'}(K^\bullet),\mathrm H^{q''}(K'^\bullet))
  \tag{11.4.6.2}
\]
(cf. M, XVII, 2).

Ici $\mathrm R^qT(K^\bullet,K'^\bullet)$ est le bicomplexe
$(\mathrm R^qT(K^i,K'^j))_{(i,j)\in\mathbb Z\times\mathbb Z}$ et le second
membre de (\hyperref[0.11.4.6.1-fr]{11.4.6.1}) est sa cohomologie quand on le
considère comme complexe simple.

En outre, la première suite spectrale est toujours régulière, et les deux
suites spectrales sont birégulières lorsqu'il existe $n$ tel que tout objet de
$\mathcal C$ et tout objet de $\mathcal C'$ admettent une résolution injective
de longueur $\leq n$, ou lorsque $K^\bullet$ et $K'^\bullet$ sont limités
inférieurement~; dans ce dernier cas, on peut en outre omettre l'hypothèse que
les limites inductives existent dans $\mathcal C$ et $\mathcal C'$.

Si $K_1^\bullet$, $K_1'^\bullet$ sont deux autres complexes de $\mathcal C$ et
$\mathcal C'$ respectivement, $f$, $g$ deux morphismes homotopes de
$K^\bullet$ dans $K_1^\bullet$, $f'$, $g'$ deux morphismes homotopes de
$K'^\bullet$ dans $K_1'^\bullet$, alors les morphismes
$\mathrm R^\bullet T(K^\bullet,K'^\bullet)\to
\mathrm R^\bullet T(K_1^\bullet,K_1'^\bullet)$ déduits de $f$ et $f'$ d'une
part, de $g$ et $g'$ d'autre part, sont identiques, et il en est de même pour
les morphismes des suites spectrales d'hypercohomologie.

On généralise aisément à un multifoncteur covariant additif quelconque.
\end{env}

\begin{proposition}[11.4.7]
\label{0.11.4.7-fr}
Supposons que pour tout objet injectif $I$ de $\mathcal C$ (resp. $I'$ de
$\mathcal C'$), $A'\mapsto T(I,A')$ (resp. $A\mapsto T(A,I')$) soit un
foncteur exact. Alors, avec les notations de
(\hyperref[0.11.4.6-fr]{11.4.6}), on a des isomorphismes canoniques
\[
  \label{0.11.4.7.1-fr}
  \mathrm R^\bullet T(K^\bullet,K'^\bullet)
  \xrightarrow{\sim}\mathrm H^\bullet(T(L^{\bullet\bullet},K'^\bullet))
  \xrightarrow{\sim}\mathrm H^\bullet(T(K^\bullet,L'^{\bullet\bullet}))
  \tag{11.4.7.1}
\]
où les deux derniers termes sont la cohomologie des complexes simples définis
par les tricomplexes $T(L^{\bullet\bullet},K'^\bullet)$ et
$T(K^\bullet,L'^{\bullet\bullet})$ respectivement.
\end{proposition}

\begin{proof}
Définissons par exemple le premier de ces isomorphismes. Le quadricomplexe
$T(L^{\bullet\bullet},L'^{\bullet\bullet})$ peut être considéré comme un
\emph{bicomplexe}, en prenant comme degrés de $T(L^{ij},L'^{hk})$ les nombres
$i+j+h$ et $k$. Comme pour chaque $h$, $L'^{h,\bullet}$ est une
\emph{résolution} de $K'^h$, on a, pour ce bicomplexe, en vertu de l'hypothèse
sur $T$,
$\mathrm H_{\mathrm{II}}^q(T(L^{\bullet\bullet},L'^{\bullet\bullet}))=0$
pour $q\neq0$ et
$\mathrm H_{\mathrm{II}}^0(T(L^{\bullet\bullet},L'^{\bullet\bullet}))
=T(L^{\bullet\bullet},K'^\bullet)$~; la première suite spectrale de ce
bicomplexe est donc \emph{dégénérée}~; comme $L'^{hk}=0$ pour $k<0$, cette
suite est en outre régulière (\hyperref[0.11.3.3-fr]{11.3.3}), et la
conclusion résulte donc de (\hyperref[0.11.1.6-fr]{11.1.6}).
\end{proof}

On a des résultats analogues pour un multifoncteur covariant en un nombre
quelconque $n$ d'arguments~: dans le calcul de l'hypercohomologie, il n'est pas
nécessaire de remplacer \emph{tous} les complexes par une résolution de
Cartan-Eilenberg, mais seulement $n-1$ d'entre eux, pourvu que, lorsqu'on fixe
$n-1$ arguments quelconques en leur donnant pour valeurs des objets
\emph{injectifs}, le foncteur covariant en l'argument restant soit
\emph{exact}.

\subsection{Passage à la limite inductive dans l'hypercohomologie}
\label{subsection:0.11.5-fr}

\begin{lemma}[11.5.1]
\label{0.11.5.1-fr}
Soit $K^\bullet=(K^i)_{i\in\mathbb Z}$ un complexe de $\mathcal C$, et pour
tout entier $r\in\mathbb Z$, soit $K_{(r)}^\bullet$ le complexe tel que
$K_{(r)}^i=0$ pour $i<r$, $K_{(r)}^i=K^i$ pour $i\geq r$. Soit $T$ un
foncteur covariant additif
